% Part of main.tex; compile main.tex.
% ---------------------------------------------------------------------------
\section{Survey Methodology}
\label{sec:survey}

In this section, we describe how we selected and coded the evaluations in
our survey, and we present our findings.

\paragraph{Scope and search.} We searched for evaluations that quantify the
benefit of asynchronous or overlapped I/O, together with the benchmarks those
evaluations are built on. Starting from the HDF5 asynchronous VOL
papers~\cite{tang2019async,tang2022async} and h5bench~\cite{li2021h5bench},
we followed citations backward and forward. In addition, we searched IEEE
Xplore, the ACM Digital Library, and Google Scholar for combinations of
``asynchronous I/O,'' ``background thread,'' ``overlap,'' ``checkpoint,'' and
``staging'' with the names of the major I/O libraries (HDF5, ADIOS and
ADIOS2, PDC, MPI-IO, and POSIX). Because benchmark tools encode their
methodology in code as much as in prose, we also inspected the released
source code of h5bench, \texttt{adios\_iotest}, and DLIO.

\paragraph{Inclusion criteria.} A source is included if it meets two
criteria. First, it reports a quantitative evaluation of asynchronous I/O, or
it is a benchmark designed to measure overlap. Second, the treatment of the
application's computation phase can be determined from its text or released
source code. We identified 26 sources, and 17 of them meet both criteria.
These 17 works are listed in Table~\ref{tab:survey}. The other nine sources
are library documentation, tutorials, and design papers whose evaluation has
no interleaved computation. They provide context in
Section~\ref{sec:related} but are not coded.

\paragraph{Coding.} For each included work, we record the I/O library and
five attributes:
\begin{enumerate}
  \item The \emph{compute phase}, which is \texttt{sleep}, a synthetic kernel
        that consumes resources (busy), a real application, unspecified, or
        none.
  \item Whether the paper \emph{describes} what the compute phase does.
  \item Whether it reports an \emph{application-side cost}, such as the
        slowdown of the compute phase or the end-to-end time, rather than
        only the observed I/O time.
  \item Whether the \emph{core allocation} for background threads or servers
        is stated.
  \item Whether the results are validated on a \emph{real application} with
        its own computation.
\end{enumerate}
Each attribute is marked as \yes{} (yes), \partly{} (partially or as a
caveat), \no{} (no), \na{} (not applicable), or ? (not determinable from the
sources we examined).

\begin{table*}[t]
\centering
\small
\setlength{\tabcolsep}{4.5pt}
\begin{tabular}{@{}l l l c c c c@{}}
\toprule
Work & Stack & Compute phase & Described & App-side cost & Cores stated & Real app \\
\midrule
h5bench~\cite{li2021h5bench}                 & HDF5            & sleep                      & \yes & \no  & ?    & \no  \\
Tang et al.\ '19~\cite{tang2019async}        & HDF5 async      & unspecified ``simulated''  & \no  & \no  & ?    & \no  \\
Tang et al.\ '22~\cite{tang2022async}        & HDF5 async      & sleep                      & \yes & \no  & \partly & \partly \\
Ravi et al.\ '23~\cite{ravi2023evaluating}   & HDF5 async      & sleep (30\,s)              & \yes & \no  & \partly & ?   \\
Tang et al.\ '18~\cite{tang2018pdc}          & PDC             & sleep                      & \yes & \no  & \yes & \no  \\
Tang et al.\ '19b~\cite{tang2019pdctune}     & PDC             & unspecified                & \no  & \no  & \yes & \partly \\
Skel~\cite{logan2011skel}                    & ADIOS           & none (I/O-only skeleton)   & \na  & \no  & \na  & \no  \\
\texttt{adios\_iotest}~\cite{adios2_iotest}  & ADIOS2          & sleep or busy loop         & \yes & \no  & \na  & \no  \\
DLIO~\cite{devarajan2021dlio,dlio_github}    & DL loaders      & sleep (GPU-profiled)       & \yes & \partly & ?   & \no  \\
Harmonia~\cite{kougkas2018harmonia}          & burst buffer    & unspecified math kernels   & \partly & ?   & ?    & \yes \\
Gossman et al.\ '21~\cite{gossman2021aggregated} & VeloC / POSIX & none (flush timed alone)  & \na  & \no  & ?    & \no  \\
\midrule
Jain et al.\ '22~\cite{jain2022flashx}       & HDF5 async      & real (Flash-X)             & \yes & \yes & ?    & \yes \\
Tseng et al.\ '21~\cite{tseng2021demystifying} & generic       & CPU/memory stressors       & \yes & \yes & ?    & \yes \\
Tarraf et al.\ '24~\cite{tarraf2024behind}   & various         & real (WaComM++, HACC-IO)   & \yes & \yes & ?    & \yes \\
IMB-IO~\cite{intel_imb_io}                   & MPI-IO          & calibrated CPU kernel      & \yes & \yes & \na  & \no  \\
COMB~\cite{lawry2002comb}                    & MPI comm.       & busy work                  & \yes & \yes & \na  & \no  \\
Denis \& Trahay~\cite{denis2016overlap}      & MPI comm.       & busy work                  & \yes & \yes & \na  & \no  \\
\bottomrule
\end{tabular}
\caption{How surveyed evaluations treat the computation phase. Upper block:
works whose compute phase cannot slow down (sleep), is unspecified, or is
absent. Lower block: works whose compute phase consumes resources.
\yes{}~yes, \partly{}~partially or as a caveat, \no{}~no, \na{}~not
applicable, ?~not determinable from the sources examined.}
\label{tab:survey}
\end{table*}

\paragraph{Findings.} Table~\ref{tab:survey} supports four observations.
\begin{enumerate}
  \item \textbf{Sleep is the most common treatment.} Of the 14 works that
        evaluate I/O (excluding the two MPI communication benchmarks and the
        generic study of Tseng et al.), 6 use \texttt{sleep}, 3 use a compute
        phase they do not describe, and 2 have no compute phase. Only 3 works
        (Jain et al., Tarraf et al., and IMB-IO) run the I/O against a
        computation that can be slowed down by it.
  \item \textbf{Sleep-based studies report I/O-side metrics.} Five of the six
        sleep-based works report only I/O time or bandwidth. DLIO reports
        end-to-end throughput, but with a sleeping compute phase it cannot
        register a slowdown. Some works state explicitly that the reported
        time ``excludes the computation time between the I/O
        operations''~\cite{tang2019async}. None of the works in the upper
        block reports the slowdown of the compute phase.
  \item \textbf{Benchmark tools make sleep the default.} h5bench,
        \texttt{adios\_iotest}, and DLIO all expose the compute phase as a
        sleep duration. Studies that build on these tools inherit this
        choice, often without stating it.
  \item \textbf{Evaluations with real computation find smaller benefits and
        identify contention as the cause.} All works in the lower block
        report application-side costs. In addition, the works that evaluate
        I/O attribute the reduced benefit to competition for CPU time,
        memory bandwidth, or the
        network~\cite{jain2022flashx,tseng2021demystifying,tarraf2024behind}.
\end{enumerate}

\summarybox{Sleep is the most common way to emulate computation in the
evaluation of asynchronous I/O, and the most popular benchmark tools use it
by default. Sleep-based studies report I/O-side metrics only, so they cannot
observe a slower computation. The few evaluations that use real
computation report smaller benefits and attribute the difference to
contention.}
